\documentclass[10pt,a4paper]{amsart}
\usepackage[margin=19mm]{geometry}
\usepackage{amsmath,amssymb,mathtools}
\usepackage[scr=rsfs,cal=euler]{mathalfa}
\usepackage{xfrac}
\usepackage[unicode,hidelinks,bookmarks=false]{hyperref}
\newcommand{\ie}{i.e.\ }
\newcommand{\cf}{cf.\ }
\newcommand{\resp}{respectively\ }
\newcommand{\Subsection}{\subsection}
\providecommand{\nobreakdash}{\nobreak\hbox{-}}
\setlength{\emergencystretch}{2em}
\multlinegap0pt
\usepackage{bm}
\usepackage{bbm}
\usepackage{dsfont}
\usepackage{xspace}
\usepackage{cancel}
\usepackage{mathtools}
\usepackage{amsthm}
\makeatletter
\@ifundefined{thm}{\newtheorem{thm}{Thm}}{}
\@ifundefined{prop}{\newtheorem{prop}[thm]{Proposition}}{}
\makeatother
\DeclareMathOperator{\Drm}{D}
\DeclareMathOperator{\dom}{dom}
\DeclareMathOperator{\grad}{grad}
\DeclareMathOperator{\id}{id}
\def\restr{r}
\newcommand\calC{\mathcal C}
\newcommand{\M}{\mathcal{M}}
\newcommand{\Tan}{\mathrm{T}}
\newcommand{\invRet}{\mathcal L}
\title[Riemannian Multilevel Optimization with Application to Const]{Riemannian Multilevel Optimization with Application to Constrained Energy Minimization Problems}
\author{Yara Elshiaty, Stefania Petra, Jonas P\"uschel, Tatjana Stykel, Ferdinand-Joseph Vanmaele}
\date{Source: arXiv:2607.09517v1 (CC BY 4.0)}
\begin{document}
\maketitle

\noindent\emph{Source and licence.} Riemannian Multilevel Optimization with Application to Constrained Energy Minimization Problems. Version arXiv:2607.09517v1 (CC BY 4.0), distributed under the Creative Commons Attribution 4.0 International licence, as recorded in the arXiv licence metadata for this exact version and shown on the abstract page https://arxiv.org/abs/2607.09517v1. Full rights notice: licenses/arxiv-2607.09517-CC-BY-4.0.txt.\par\smallskip

\noindent\textit{Editorial scope.} Counted unit: the Prop whose source label is \texttt{prop:algebraic1} in the article (source0.tex lines 672--674, source label \texttt{prop:algebraic1}), delivered together with the article's own complete original proof of it (source0.tex lines 676--697, one \texttt{proof} environment of 22 lines). No mathematics is written, repaired or re-derived: statement and proof are the article's own text line for line. Required context retained uncounted: eq:riem\_coarse\_model (lines 475--478); eq:r\_algebraic\_prol\_restr (lines 662--665). Every definition, display and label of those lines is retained unchanged. Omitted from this package and not redistributed: 1--474, 479--661, 666--671, 675--675, 698--2035. Nothing inside a retained excerpt is omitted or altered: every retained body is delivered exactly as the article's lines are written, so no omission receipt of any kind accompanies this package. Citations: the retained excerpts carry no citation command at all, so no bibliography entry of the article is transcribed and no reference is redistributed. Reference adaptation: none was needed and none was made; every reference inside the delivered unit resolves inside it, so no unresolved reference or citation is produced. Printed slips kept literal: no printed word, symbol, hypothesis or number of the counted statement or of its proof was altered. Layout and class: the article's own class is \texttt{amsart} with packages including appendix, svg, graphicx, url, amsmath, amssymb, amsthm, bm, bbm, booktabs, polynom, contour, xcolor, pgfplots; the wrapper is the lane's pinned \texttt{amsart} editorial wrapper at 10pt on A4, which restates the theorem-like environments the retained text needs and adds the article's own macro definitions that the retained text needs (\texttt{\textbackslash Drm}, \texttt{\textbackslash M}, \texttt{\textbackslash Tan}, \texttt{\textbackslash calC}, \texttt{\textbackslash dom}, \texttt{\textbackslash grad}, \texttt{\textbackslash id}, \texttt{\textbackslash invRet}, \texttt{\textbackslash restr}), with extra wrapper packages (bm, bbm, dsfont, xspace, cancel, mathtools). The delivered numbering is the unit's own. Assets: no retained span contains a figure environment, a graphics-inclusion command, a PDF-page inclusion, a table or a TikZ picture; no image file is copied into this package, the package declares no asset dependency and no image file is redistributed with it. Licence: version arXiv:2607.09517v1 of this article is distributed under the Creative Commons Attribution 4.0 International licence, as recorded in the arXiv licence metadata for this exact version and shown on the abstract page https://arxiv.org/abs/2607.09517v1. A full rights notice is delivered at licenses/arxiv-2607.09517-CC-BY-4.0.txt. Characters: every retained excerpt is pure ASCII, so the delivered engine, XeLaTeX with its default Latin Modern text font, reports no missing character.
	\begin{equation} \label{eq:riem_coarse_model}
		\min\limits_{z \in \calC_{y}} \; q(z) = f_H(z) - \langle w, \invRet^H_{y}(z)\rangle_{y}
		\enskip \text{with }\; w = \grad f_H(y) - R_{x}^{y}(\grad f_h(x)),
	\end{equation}

	\begin{equation} \label{eq:r_algebraic_prol_restr}
		P_y^x = \Drm \restr(x)^-, \qquad\quad
		R_x^y = (P_y^x)^\ast,    
	\end{equation}

	\begin{prop}\label{prop:algebraic1} Let $x \in \dom(r)$ and $y = \restr(x)$ for a~point restriction $\restr$.
		If $f_h = f_H \circ \restr$ on $\dom(r)$, then for the vector transfer operators \eqref{eq:r_algebraic_prol_restr}, the coarse model objective in~\eqref{eq:riem_coarse_model} satisfies $q(z) = f_H(z)$ for all $z\in \calC_{y}$.
	\end{prop}

	\begin{proof}
		Since $\Drm \restr(x)$ is surjective, its right inverse $P_{y}^{x}$ is injective, and $R_{x}^{y} = (P_{y}^{x})^\ast$ is surjective. 
		% This is correct but not used in the proof
		By the definition of the right inverse, we obtain that
		\[
		R_{x}^{y} \circ \Drm \restr(x)^\ast 
		= \big(\Drm \restr(x)\circ\Drm \restr(x)^-\big)^\ast
		= \id_{\Tan_{y}\M_H}.
		\]
		By chain rule, we further have 
		\begin{align*}
			\grad f_h(x) &= \grad (f_H \circ r)(x) 
			= \Drm \restr(x)^\ast(\grad f_H(\restr(x))).
		\end{align*}
		Therefore, 
		\[
		w 
		= \grad f_H(y)-\big(R_{x}^{y}\circ \Drm \restr(x)^\ast\big)(\grad f_H(y))
		= 0,
		\]
		and consequently $q(z) = f_H(z)$  for all $z \in \calC_{y}$.
	\end{proof}

\end{document}
